\begin{lemma}
For every real \(\varepsilon>0\), there is an integer \(n_0\) such that for
every integer \(n\ge n_0\) there exist an integer \(m\), a real \(p\), and a
sample
\[
\omega\in\noderef{TwoBiteSample}(n,m,p)
\]
such that
\[
m=\noderef{TwoBiteNaturalReducedVertexCount}(n),\qquad
p=\noderef{TwoBiteEdgeProbability}\!\left(\frac12,n\right),
\]
\[
0<\noderef{TwoBiteSampleWeight}(\omega),
\]
the shadow graph \(\noderef{TwoBiteFinalGraph}(\omega)_0\) has no triangle,
and
\[
\noderef{IndependenceNumberLess}\!\left(
\noderef{TwoBiteFinalGraph}(\omega)_0,
(1+\varepsilon)\sqrt{(n:\mathbb R)\log(n:\mathbb R)}\right).
\]
\end{lemma}
\begin{proof}
Fix \(\varepsilon>0\).  Apply
\(\noderef{ParameterHierarchyChoice}\) to choose
\(\eta\), \(\varepsilon_1\), \(\varepsilon_2\), and a hierarchy threshold
\(n_{\mathrm h}\) such that
\[
0<\eta<\varepsilon
\qquad\text{and}\qquad
\noderef{ParameterHierarchy}(\eta,\varepsilon_1,\varepsilon_2,n_{\mathrm h})
\]
holds.  Put \(\beta=\frac12\),
\[
M(n)=\noderef{TwoBiteNaturalReducedVertexCount}(n),\qquad
p(n)=\noderef{TwoBiteEdgeProbability}\!\left(\frac12,n\right).
\]
By \(\noderef{TwoBiteParameterRelations}\) and the definition of
\(\noderef{TwoBiteEdgeProbability}\), after increasing the final threshold if
necessary,
\[
p(n)=\frac12\sqrt{\frac{\log n}{n}}
\]
and \(p(n)>0\).  The integer \(M(n)\) is the floor of the paper's real base size
\(n/\log^2 n\), so it is the constructible base-graph size used by the finite
probability space.

Apply \(\noderef{NoLargeIndependentSetWhp}\) with target \(\varepsilon\),
inner hierarchy parameter \(\eta\), and \(\beta=\frac12\).  Its conclusion says,
in the sense of \(\noderef{WithHighProbability}\), that the probabilities
\[
P_n=\noderef{TwoBiteEventProbability}\bigl(n,M(n),p(n),\,
\{\omega:\mathrm{IndependenceNumberLess}
(G(\omega),(1+\varepsilon)\sqrt{n\log n})\}\bigr)
\]
tend to \(1\).  Hence, after increasing the threshold again, \(P_n>0\) for every
\(n\) under consideration.

Fix such an \(n\).  By the definition
\(\noderef{TwoBiteEventProbability}\), the number \(P_n\) is the finite sum of
\(\noderef{TwoBiteSampleWeight}(\omega)\) over exactly those two-bite samples
\(\omega\) for which
\[
\mathrm{IndependenceNumberLess}
(G(\omega),(1+\varepsilon)\sqrt{n\log n})
\]
holds.  If every summand in this finite sum were \(0\), the sum would be \(0\),
contrary to \(P_n>0\).  Choose a summand with positive weight, and call the
corresponding sample \(\omega\).  Then
\[
0<\noderef{TwoBiteSampleWeight}(\omega)
\]
and the final graph \(G(\omega)\) has independence number less than
\((1+\varepsilon)\sqrt{n\log n}\), which is precisely the predicate
\(\noderef{IndependenceNumberLess}\) appearing in the statement.

Finally, \(\noderef{TwoBiteDeletionTriangleFree}\) applies to every two-bite
sample, independently of the probabilistic estimates.  Applying it to the
chosen \(\omega\) shows that the simple final graph \(G(\omega)\) has no
triangle.  Thus the chosen parameters \(m=M(n)\), \(p=p(n)\), and the
positive-weight sample \(\omega\) satisfy all required conclusions:
the rounded definition of \(m\), the displayed formula for \(p\),
triangle-freeness of the final graph, and the required independence-number
bound.
\end{proof}
